% 2013 年辽宁卷理科第 8 题：程序框图。
% 来源 PDF：普通高考/2013/2013辽宁理.pdf，第 1 页（印面 821）中栏。
% 流程：开始 → 输入 n → S=0,i=2 → i≤n? 是：S=S+1/(i^2−1)，i=i+2，回到判断；
% 否：输出 S → 结束。
\begin{tikzpicture}[font=\footnotesize, >=Stealth, line width=0.45pt,
  node distance=5mm,
  term/.style={draw, rounded corners=6pt, minimum height=1.8em, inner xsep=9pt},
  io/.style={draw, trapezium, trapezium left angle=70, trapezium right angle=110,
             minimum height=1.8em, inner xsep=6pt},
  box/.style={draw, minimum height=1.9em, inner xsep=6pt},
  dec/.style={draw, diamond, aspect=2.2, inner sep=1pt}]
  \node[term] (start) {开始};
  \node[io, below=5mm of start] (input) {输入 $n$};
  \node[box, below=5mm of input] (init) {$S=0$，$i=2$};
  \node[dec, below=6mm of init] (cond) {$i\leqslant n$};
  \node[box, below=7mm of cond] (acc) {$S=S+\dfrac{1}{i^2-1}$};
  \node[box, below=5mm of acc] (inc) {$i=i+2$};
  \node[io, right=15mm of cond] (output) {输出 $S$};
  \node[term, below=6mm of output] (end) {结束};
  \draw[->] (start) -- (input);
  \draw[->] (input) -- (init);
  \draw[->] (init) -- (cond);
  \draw[->] (cond) -- node[right] {是} (acc);
  \draw[->] (acc) -- (inc);
  \draw[->] (cond) -- node[above] {否} (output);
  \draw[->] (output) -- (end);
  % 循环支路：从 i=i+2 左侧引出，向上回到判断框入口
  \draw[->] (inc.west) -- ++(-1.45,0) |- (cond.west);
\end{tikzpicture}
